\documentclass[11pt]{article}
\usepackage[margin=2.4cm]{geometry}
\usepackage{amsmath,amssymb}
\usepackage{booktabs}
\usepackage{graphicx}
\usepackage{multirow}
\usepackage{caption}
\usepackage[hidelinks]{hyperref}
\captionsetup{font=small,labelfont=bf}
\newcommand{\fcntm}{FC-NTM}
\renewcommand{\thetable}{S\arabic{table}}
\renewcommand{\thefigure}{S\arabic{figure}}
\title{\bfseries Supplementary Material\\[4pt]
\large Auxiliary Experiments, Controls and Metric Definitions}
\author{}
\date{}
\begin{document}
\maketitle

\section{Auxiliary experiments and controls}

\subsection{Long-document truncation control}

The GitHub background field has a long tail (up to 20532 tokens after preprocessing) while the
sentence-embedding baselines truncate at 128 tokens. To check whether the ordering of the main
tables is caused by this truncation, the three baselines were re-run with window pooling: each
document is split into 128-token windows with the model's own tokenizer, the windows are encoded
separately and their vectors averaged, so the embeddings cover the full document; all other
settings are unchanged. Table~\ref{tab:chunked} reports the result.

\begin{table}[t]
\centering
\caption{Long-document control for the sentence-embedding baselines ($C_V$, seeds 11/23/37, mean$\pm$standard deviation). ``Truncated'' encodes only the first 128 tokens (the setting used in the main tables); ``window pooling'' splits a document into 128-token windows with the model's own tokenizer and averages the window vectors, so every token enters exactly one window. Top2Vec uses Doc2Vec document vectors and does not accept window pooling.}
\label{tab:chunked}
\footnotesize
\setlength{\tabcolsep}{3pt}
\begin{tabular}{@{}lcccc@{}}
\toprule
model & GitHub Truncated & GitHub Window pooling & AG News Truncated & AG News Window pooling \\
\midrule
BERTopic & 0.540$\pm$0.015 & 0.582$\pm$0.009 & 0.488$\pm$0.022 & 0.518$\pm$0.011 \\
Top2Vec & 0.493$\pm$0.009 & n/a & 0.498$\pm$0.032 & n/a \\
CombinedTM & 0.643$\pm$0.022 & 0.656$\pm$0.018 & 0.526$\pm$0.037 & 0.488$\pm$0.030 \\
\bottomrule
\end{tabular}
\end{table}

Window pooling raises BERTopic from 0.540 to 0.582 and CombinedTM from 0.643 to 0.656 on GitHub,
and BERTopic from 0.488 to 0.518 while CombinedTM moves from 0.526 to 0.488 on AG News. Even at
their best (GitHub CombinedTM 0.656, AG News BERTopic 0.518) the pooled baselines remain below the
$\beta$ read-out of \fcntm{} (0.701 and 0.586) and well below its class TF--IDF read-out (0.727 and
0.612), so the ordering of bag-of-words models and sentence-embedding baselines is not an artefact
of truncation.

\subsection{Within-corpus length-ratio ladder}

The cross-corpus comparison mixes the length ratio with domain and genre. To separate part of this
confound, the civil corpus is truncated to the first 60 and 30 words of every abstract (length
ratio $11.29\to3.99\to2.02$) and the full, no-cross and no-mask configurations are re-run for the
same 6000/2000 record IDs and the same 7685-word vocabulary (700-epoch budget, all other settings
unchanged). Table~\ref{tab:ladder} reports the result; three observations follow. First, the
$\beta$ read-out of the full model decreases monotonically with truncation
($0.633\to0.595\to0.534$) while the class TF--IDF read-out does not
($0.636\to0.645\to0.606$): the effect of truncation depends on the read-out. Second, the paired
difference of removing cross-field reconstruction does not change monotonically with the length
ratio ($\beta$: $-0.022\to+0.006\to-0.006$; class TF--IDF: $-0.019\to+0.001\to-0.023$). Third, the
cost of removing word masking also fails to follow the ratio ($\beta$: $-0.010\to-0.002\to-0.027$;
class TF--IDF: $-0.011\to-0.017\to-0.025$), and the read-out with the smallest cost differs. Under
a controlled change of the abstract length alone, component importance is therefore neither
monotone in the ratio nor consistent across read-outs, which agrees with the cross-corpus result
that the length ratio is a description rather than a pre-training criterion.

\begin{table}[t]
\centering
\caption{Within-corpus length-ratio ladder on the civil corpus ($K=30$, seeds 11/23/37). $\Delta C_V$ is the paired difference of removing one component relative to the full model (negative means the component matters); parentheses count positive seeds. The length ratio is the mean token ratio of abstracts to titles over the same records and vocabulary.}
\label{tab:ladder}
\small
\setlength{\tabcolsep}{4pt}
\begin{tabular}{@{}llccc@{}}
\toprule
Variant(Length ratio) & read-out & Full $C_V$ & $-\hspace{2pt}$Cross-field recon. $\Delta C_V$
 & $-\hspace{2pt}$Word masking $\Delta C_V$ \\
\midrule
Original(11.29) & $\beta$ & 0.633 & $-$0.022 (0/3) & $-$0.010 (1/3) \\
Original(11.29) & Class TF--IDF & 0.636 & $-$0.019 (0/3) & $-$0.011 (0/3) \\
abs60(3.99) & $\beta$ & 0.595 & $+$0.006 (3/3) & $-$0.002 (2/3) \\
abs60(3.99) & Class TF--IDF & 0.645 & $+$0.001 (1/3) & $-$0.017 (0/3) \\
abs30(2.02) & $\beta$ & 0.534 & $-$0.006 (1/3) & $-$0.027 (0/3) \\
abs30(2.02) & Class TF--IDF & 0.606 & $-$0.023 (0/3) & $-$0.025 (0/3) \\
\bottomrule
\end{tabular}
\end{table}

\subsection{Controls for the occlusion protocol}

Three controls test what the relative degradation can show: (i) a constant unigram model that only
uses training term frequencies and reads no test document; (ii) a shuffled summary, where the
summary of each document is replaced by that of another document while the model is unchanged; and
(iii) the same architecture trained without whole-field dropout. All three use the same masks,
seeds and per-token NLL. Tables~\ref{tab:occl-controls} and~\ref{tab:fielddrop-control} report the
results.

\begin{table}[t]
\centering
\caption{Mechanism controls for the partial-occlusion protocol (seeds 11/23/37). Each cell shows the mean$\pm$standard deviation of the relative degradation $\mathrm{NLL}(100\%)/\mathrm{NLL}(25\%)$ and the absolute NLL at full occlusion. The constant unigram model reads no test input; ``shuffled summary'' replaces each summary by that of another document; ``true summary'' is the main model.}
\label{tab:occl-controls}
\small
\setlength{\tabcolsep}{4pt}
\begin{tabular}{@{}lccc@{}}
\toprule
corpus & constant unigram & shuffled summary(\fcntm{}) & \fcntm{}(true summary) \\
\midrule
GitHub & 0.999$\pm$0.001 / 7.382 & 1.091$\pm$0.002 / 8.567 & 1.029$\pm$0.003 / 7.961 \\
AG News & 1.001$\pm$0.000 / 7.809 & 1.096$\pm$0.001 / 8.622 & 1.018$\pm$0.001 / 7.828 \\
civil corpus & 1.000$\pm$0.001 / 7.875 & 1.084$\pm$0.001 / 8.997 & 1.010$\pm$0.001 / 8.160 \\
\bottomrule
\end{tabular}
\end{table}

\begin{table}[t]
\centering
\caption{Effect of training-time field dropout on the partial-occlusion curve (relative degradation, seeds 11/23/37). The two columns are the same architecture with and without whole-field dropout during training.}
\label{tab:fielddrop-control}
\small
\setlength{\tabcolsep}{4pt}
\begin{tabular}{@{}lcc@{}}
\toprule
corpus & dropout(mainmodel) & No dropout \\
\midrule
GitHub & 1.029$\pm$0.003 & 1.028$\pm$0.007 \\
AG News & 1.018$\pm$0.001 & 1.020$\pm$0.003 \\
civil corpus & 1.010$\pm$0.001 & 1.012$\pm$0.001 \\
\bottomrule
\end{tabular}
\end{table}

Three conclusions follow. First, the constant unigram model has a relative degradation of almost
one on all three corpora (0.999--1.001) and its absolute NLL at full occlusion is lower than that
of \fcntm{} on GitHub and the civil corpus (7.382 versus 7.961 and 7.875 versus 8.160) and close on
AG News (7.809 versus 7.828): a model that reads no test input can be ``stable'' under this ratio,
and the cross-model comparison of absolute likelihood mixes in the calibration of the word
distributions. The ratio is therefore a diagnostic under a specified baseline set, not evidence of
field complementarity by itself. Second, replacing the summary by that of another document raises
the relative degradation from 1.029/1.018/1.010 to 1.091/1.096/1.084 and worsens the absolute NLL
(8.567/8.622/8.997 versus 7.961/7.828/8.160): the model really uses the \emph{content} of the
summary field, which supports ``the summary participates in inference'' more directly than the
ratio itself. Third, removing training-time dropout leaves the ratio almost unchanged
(1.028/1.020/1.012 versus 1.029/1.018/1.010). A same-budget single-encoder control is not
compared under this occlusion protocol, so the behaviour on these curves cannot be attributed
to the two-stream structure alone.

\subsection{Sensitivity of the $\beta$ read-out}

The main tables read $\beta$ from the decoder parameter rows $B_k$, while the document generative
distribution of \fcntm{} is $\operatorname{softmax}(\operatorname{BN}(\theta B))$; the two are not
equivalent. Feeding a unit topic input $e_k$ into the decoder gives the generative distribution,
whose top-10 words overlap the raw read-out with an average Jaccard of only 0.36--0.42 and whose
first word differs for 56\%--74\% of the topics; the $C_V$ and NPMI of the same models also drop
from the raw values (GitHub 0.701$\to$0.625 and 0.233$\to$0.192; AG News 0.586$\to$0.541 and
0.137$\to$0.100; civil corpus 0.633$\to$0.584 and 0.154$\to$0.118), as Table~\ref{tab:beta-readout}
shows. A translation test makes the same point: adding one shared per-word vector to all topic
rows and shifting the BatchNorm running mean leaves the generative distribution pointwise unchanged
(maximum probability difference $\leq5\times10^{-7}$) while 100\% of the raw top-10 lists change.
The class TF--IDF read-out is built from $\theta$ and term counts and is not affected; both
$\beta$ readings are reported, and the generative reading is the one that reflects the model's
generative semantics.

\begin{table}[t]
\centering
\caption{Sensitivity of the $\beta$ read-out (main checkpoints of the three corpora, means over seeds 11/23/37). ``raw'' reads the decoder parameter rows $B_k$; ``BN'' reads the generative distribution $\operatorname{softmax}(\operatorname{BN}(B_k))$ for a unit topic input; Jaccard is the average per-topic overlap of the top-10 words and the top-1 change rate the share of topics whose first word differs.}
\label{tab:beta-readout}
\small
\setlength{\tabcolsep}{4pt}
\begin{tabular}{@{}lcccccc@{}}
\toprule
corpus & raw $C_V$ & BN $C_V$ & raw NPMI & BN NPMI & top-10 Jaccard & top-1 change rate \\
\midrule
GitHub & 0.701 & 0.625 & 0.233 & 0.192 & 0.369 & 0.74 \\
AG News & 0.586 & 0.541 & 0.137 & 0.100 & 0.360 & 0.62 \\
civil corpus & 0.633 & 0.584 & 0.154 & 0.118 & 0.421 & 0.56 \\
\bottomrule
\end{tabular}
\end{table}

\section{Data and metric notes}

\subsection*{Training details of the baselines}

The four baseline families use the following settings. \emph{Gibbs sampling}: BTM is a collapsed
Gibbs sampler over biterms (word pairs within a window of 12 tokens, at most 150 sampled pairs per
document, $\alpha=50/K$, $\beta=0.01$, 300 full sweeps); LDA is fitted by gensim's batch variational
Bayes (50 passes, 120 iterations per pass, symmetric $\alpha$). \emph{Matrix factorisation}: NMF
minimises the Kullback--Leibler divergence between the count matrix and $WH$ with multiplicative
updates (600 iterations, tolerance $10^{-4}$). \emph{Neural variational models}: ProdLDA, ETM,
CTM-lite and FC-NTM share the Adam loop, mini-batch size and budget (700 epochs, 20\% held-out early
stopping); the baselines use learning rate $2\times10^{-3}$ and the main FC-NTM configuration
$5\times10^{-4}$ selected on development data, and the official CombinedTM keeps its own loop
(Adam, learning rate $2\times10^{-3}$, batch size 64, 700 epochs, early stopping on its own held-out
split). \emph{Embedding clustering}: BERTopic runs UMAP followed by HDBSCAN with fixed random
seeds, and Top2Vec fits Doc2Vec document vectors followed by HDBSCAN; both choose their own topic
numbers and are reported for reference only.

\subsection{Data notes}

After stop-word removal and vocabulary filtering, a small number of identical token sequences occur
in both the training and the test split (GitHub 10 records, AG News 8 records, i.e.\ 0.50\% and
0.11\% of the test sets); they come from different original records and are a degeneration of the
filtering rather than duplicated raw records, and the proportion is too small to change the
conclusions. The GitHub background field has a long tail (8501 tokens in the training set and
20532 in the test set after preprocessing) that the sentence-embedding baselines truncate. A few
records have an empty summary after preprocessing (GitHub 6 training and 1 test record; AG News 69
training and 98 test records, up to 1.29\% of the test set); these records are kept in the main
tables with an empty summary input, for which the model has no explicit fallback branch -- the
summary encoder still runs and produces a training-parameter-dependent posterior before gating
with the background posterior (for the seed-11 checkpoint the maximum difference from a
background-only posterior is 0.115 for the GitHub record and 0.261 for the AG News records, with
the most probable topic unchanged in all but three cases). The dropout training and the occlusion
experiments therefore validate \emph{background} absence and should not be extrapolated to missing
summaries. Design iteration used 400 test records (a 1000/400 subsample per corpus); the main text
reports both the full test split and a development-free subset (1600 GitHub and 7200 AG News
records) obtained by exact record-level matching, on which the coherence is recomputed for every
stored model. That re-scoring removes direct overlap only and cannot undo the model and
configuration choices already made, so it is a sensitivity analysis rather than an untouched
confirmation set; an independent confirmation would require repositories or corpora not involved
in the development.

\subsection{Supplementary definitions of the metrics}

$\mathrm{EXCL}_{\mathrm{KADJ}}$ is a $K$-corrected variant of the Bischof--Airoldi exclusivity: the
usual $\mathrm{EXCL}$ is the mean column-normalised mass of the top ten words of each topic, and it
is mapped to $\mathrm{EXCL}_{\mathrm{KADJ}}=(\mathrm{EXCL}-1/K)/(1-1/K)$ so that random topics
score about zero and fully exclusive topics about one. This is our definition rather than the
original measure, and it does not remove the comparability issues between different $K$ or between
different score matrices, so the table only compares models with the same $K=30$. For the
\fcntm{}-cTFIDF rows the value is computed on the class TF--IDF score matrix after row
normalisation, which is not a probability distribution, so those rows are only ranked within the
same read-out and are not compared with $\beta$ rows. IRBO uses the finite-length extrapolated
form, $\mathrm{RBO}_{\mathrm{ext}}=(1-p)\sum_{d=1}^{L}p^{d-1}A_d+p^{L}A_L$ with $p=0.9$, common
depth $L$ and set overlap $A_d$ at depth $d$, and topic diversity is
$\mathrm{IRBO}=1-\overline{\mathrm{RBO}}_{\mathrm{ext}}$. Without the tail term even two identical
ten-word topic lists would obtain a ``similarity'' of $1-p^{L}=0.651$ (diversity 0.349); the bias
varies with the overlap at the final depth (it equals $p^{L}$ for identical and zero for disjoint
lists), so it is not a constant offset shared by all models.

\section{Additional tables and figures}

The table below lists the example topics referenced in the main text.

\begin{table}[t]
\centering
\caption{Example topics on GitHub issues (\fcntm{}, seed 11, $K=30$, best epoch 547 -- the same model as in the main table and figure). The words are generated from the topic--word matrix: the class TF--IDF read-out is computed first and then truncated to the summary and background vocabularies, so no word lies outside the vocabulary; topic names are assigned manually.}
\label{tab:topics}
\footnotesize
\setlength{\tabcolsep}{4pt}
\begin{tabular}{@{}p{2.7cm}p{5.9cm}p{5.9cm}@{}}
\toprule
topics(manual) & Summary-side words(top 10 ) & Background-side words(top 10 ) \\
\midrule
% github seed=11 K=30 best_epoch=547 (curated example topics, manual names)
% topic name & summary-side words (top 10) & background-side words (top 10)
Type systems and language services & type typescript function string types regression error property assignable array & type typescript function string types regression error property assignable array \\
DuckDB build and distribution & duckdb all python data sets distribution architecture linux required nightly & duckdb all python data reproducing sets distribution architecture linux required \\
Redis replication and clustering & master cluster redis node sentinel slave nodes docker failover command & master cluster redis node sentinel slave nodes docker failover command \\
Configuration and command semantics & redis command key set value some keys error one cluster & redis command key set value some keys error one cluster \\
Language feature proposals & javascript type typescript suggestion types existing runtime change feature syntax & javascript type typescript suggestion types wouldn existing runtime change feature \\
Event loop and cache metrics & idle addr sub flags cmd age multi free events name & qbuf psub oll obl idle addr sub flags cmd omem

\\   % the row terminator lives here, not in the \input file: a trailing \\
     % at the end of an \input-ed file breaks \bottomrule (Misplaced \noalign)
\bottomrule
\end{tabular}
\end{table}

\subsection{Civil topic case studies}

The table below lists four readable topics of the seed-11 civil model with one representative
record each, plus two failure cases found by scanning the labelled test records. In the first
failure case a seismicity paper is routed by its title alone to the BIM topic, and the abstract
restores the correct label; in the second, an entry whose title mentions corrosion is routed by
both settings to the related ``concrete reinforcement materials'' topic instead of the corrosion
topic. Routing uses the same majority-vote mapping protocol as the main archival task. The
representative records of the four topics are selected by topic mass on the training split and
are not held-out evidence; the two failure cases are test records. The mixed laboratory-materials
topic is an interpretability problem (mixed vocabularies), not a routing failure --- its
representative record is routed correctly.

\begin{table}[t]
\centering
\caption{Case studies on the civil corpus (\fcntm{}, seed 11, $\beta$ read-out; ``title-only'' and ``full-field'' give the routing decision of the same mapping protocol, and ``miss'' marks a decision different from the OpenAlex label; OpenAlex Work IDs are given in parentheses).}
\label{tab:civil-cases}
\scriptsize
\setlength{\tabcolsep}{3pt}
\begin{tabular}{@{}p{2.6cm}p{4.0cm}p{3.6cm}p{2.4cm}cc@{}}
\toprule
Topic(manual) & Top words($\beta$ read-out) & Representative record & OpenAlex label & title-only & full-field \\
\midrule
% civil seed=11 beta read-out (case studies, manual names; reconstructed from supplementary.pdf Table S7)
% topic name & top words (beta read-out) & representative record & OpenAlex label & title-only & full-field
Asphalt mixture design and performance & asphalt, mixture, binder, mixtures, temperature, test, performance, modified, stability, low & Synergetic Effect of NanoZnO and Trinidad Lake Asphalt for Antiaging Properties of SBS Modified Asphalt (W3007633598) & Asphalt Pavement Performance Evaluation & correct & correct \\
Seismic response of tunnels and piles & response, ground, motion, seismic, tunnel, earthquake, acceleration, foundation, pile, wave & Earthquake ground-motion parameters at stiff sites in Addis Ababa, Ethiopia (W2033407901) & Seismic Performance and Analysis & correct & correct \\
BIM and construction management & construction, information, bim, building, project, industry, management, projects, implementation, process & Exploring the Adoption of Building Information Modeling (BIM) in the Jordanian Construction Industry (W2624290857) & BIM and Construction Integration & correct & correct \\
Structural mechanics: shear, buckling, composite beams & load, capacity, finite, element, shear, beam, plate, composite, behavior, buckling & Out-of-Plane Secondary Bifurcation Buckling Behavior of Elastic Circle Pipe Arch (W2063862378) & Structural Load Bearing Analysis & correct & correct \\
Mixed laboratory materials topic (failure case) & diamond, wood, metal, cell, density, wear, segments, adsorption, matrix, saw & Pseudotsuga menziesii (Mirb.) Franco wood basic density in the Patagonia Argentina (W1511725581) & Wood Treatment and Properties & correct & correct \\
Failure case: title only misrouted & construction, information, bim, building, project, industry, management, projects, implementation, process & United States building code approach to variations in regional seismicity (W2525447080) & Seismic Performance and Analysis & miss & correct \\
Failure case: title only misrouted & strength, fiber, steel, concrete, mechanical, reinforced, fibers, specimens, compressive, tensile & Corrosion Behavior of Reinforcing steel in Concrete (W4387417777) & Concrete Corrosion and Durability & miss & miss

\\   % the row terminator lives here, not in the \input file
\bottomrule
\end{tabular}
\end{table}

\end{document}
